随着科研产出的指数级增长，如何从浩如烟海的文献中快速定位相关研究，成为科研工作者的核心痛点之一。引文网络（citation network）作为连接学术成果的结构化图谱，蕴含了丰富的“知识流向”信息：一篇论文的参考文献，本质上刻画了它在学术谱系中的来路。若能利用已有引文数据，预测一篇新论文最可能引用哪些文献，便能显著提升文献推荐与学术检索的体验。

传统方法主要依赖\emph{内容相似度}（如 TF-IDF + 余弦相似度）或\emph{协同过滤}（基于共同引用模式）\cite{kipf2017}。然而，前者难以捕捉深层语义，后者则面临严重的冷启动问题。近年来，图神经网络（GNN）在建模非欧结构数据上展现出强大能力\cite{wang2019}，为引文预测提供了新的建模工具；与此同时，以 Transformer 为代表的预训练语言模型\cite{vaswani2017,devlin2019} 能够为中文文本提供高质量的语义表示。

本文的主要贡献如下：

\begin{enumerate}
  \item 提出 Papex-GNN，一个面向中文文献的端到端引用预测框架，首次将预训练语义编码与异构图注意力结构建模统一在同一框架下；
  \item 设计了同时建模“主题相似边”与“作者协作边”的异构消息传递机制，显式区分两类结构信号；
  \item 在自建中文论文数据集上取得显著优于基线的效果，并提供了充分的消融分析。
\end{enumerate}
